\documentclass[11pt,a4paper]{article}

\usepackage[margin=1in]{geometry}
\usepackage[T1]{fontenc}
\PassOptionsToPackage{hyphens}{url}
\usepackage{hyperref}
\usepackage{listings}
\usepackage{xcolor}
\usepackage{booktabs}
\usepackage{enumitem}
\usepackage{microtype}

\emergencystretch=3em
\newcommand{\pth}[1]{\path{#1}}

\hypersetup{
    colorlinks=true,
    linkcolor=blue,
    urlcolor=blue,
    citecolor=blue,
    breaklinks=true
}

\lstset{
    basicstyle=\ttfamily\footnotesize,
    breaklines=true,
    frame=single,
    numbers=none,
    showstringspaces=false,
    keywordstyle=\color{blue},
    commentstyle=\color{gray},
    backgroundcolor=\color{gray!10}
}

\title{Closing Two Fast Commit Coverage Gaps in ext4}

\author{
    Suyamoon Pathak (241110091) \\
    CS614 Linux Kernel Programming \\
    Indian Institute of Technology Kanpur
}

\date{April 2026}

\begin{document}

\maketitle

\begin{abstract}
ext4's fast-commit feature gives \texttt{fsync()} a low-latency
commit path, but falls back to a full JBD2 commit on operations it
does not know how to log. We close two such fallbacks on
Linux~6.1.4: a 108-line patch that adds fast-commit support for
inline extended attributes, and a 46-line patch that adds it for
the COLLAPSE\_RANGE and INSERT\_RANGE modes of \texttt{fallocate}.
On the VM the xattr patch reduces JBD2 full commits on a 5000-op
\texttt{setxattr} workload by 64$\times$ (5000 to 78) and cuts
wall-time by 27\%; the fallocate patch reduces full commits on a
1000-op workload by 62.5$\times$ (1001 to 16) and cuts wall-time
by 23--26\%. Bare-metal runs reproduce the commit-count reductions
byte-identically and roughly double the wall-time gain (48\% for
xattr, 43--44\% for fallocate). Six custom crash-recovery tests
and an xfstests subset confirm no regression. We also report two
earlier candidate patches that produced no measurable improvement
and the lessons drawn from them.
\end{abstract}

\section{Introduction}

The project's goal is to optimize the consistency mechanism of
ext4 on Linux~6.1.4. The mechanism in question is JBD2 (the
journal that keeps the filesystem crash-consistent) and its
fast-commit extension introduced in Linux~5.10. Any change must
preserve crash-recovery semantics and must be validated by
measurement.

We attempted four candidate optimizations. Two of them (Section~3)
were correctly implemented patches against developer TODO comments
in the kernel but optimized paths that were hit too rarely to
measure. The two that worked (Sections~4 and~5) close fast-commit
fallbacks that fire on 100\% of a specific operation class:
\texttt{setxattr} for Candidate~3, fallocate's COLLAPSE\_RANGE and
INSERT\_RANGE for Candidate~4. Both are small (108 and 46 lines),
both reuse existing fast-commit primitives, and both reproduce
their architectural claim (a $\sim$62--64$\times$ drop in full
commits) on bare-metal hardware as well as the VM.

\paragraph{What we built.}
\begin{itemize}[leftmargin=*]
  \item Two patches against pristine 6.1.4:
        \pth{fc-inline-xattr.patch} (C3) and
        \pth{fc-fallocate-range.patch} (C4).
  \item Two benchmark scripts (xattr, fallocate) and a
        characterization harness used to pick the C4 target.
  \item Six custom crash-recovery scripts (three for each patch).
  \item Two earlier patches with postmortems explaining why they
        produced no measurable improvement.
\end{itemize}

\section{Background}

The relevant background, restricted to what the rest of the
report relies on, is the fast-commit fallback. Fast commit keeps a
small reserved region of the journal and writes a compact record
per modified inode instead of running a full JBD2 commit. When
ext4 hits an operation it cannot log this way, it calls
\texttt{ext4\_fc\_mark\_ineligible(\_,REASON,\_)} and falls back
to the full commit path. The original FastCommit
paper~\cite{atc24-fc} lists ten such reasons, including
\texttt{XATTR}, \texttt{FALLOC\_RANGE}, \texttt{CROSS\_RENAME},
and \texttt{RENAME\_DIR}. Two of those reasons
(\texttt{XATTR}~for the inline case, and \texttt{FALLOC\_RANGE}
for collapse and insert range) are what our two successful patches
remove.

\section{Candidates that did not measure (C1, C2)}

Both candidates were implementations of long-standing TODO
comments in the kernel. Both were correct (built, booted,
mounted, no dmesg errors), but neither produced a measurable
speedup on our workloads.

\subsection{C1: fast-commit barrier deferral}

The TODO sits at \texttt{fs/jbd2/commit.c} around line~454:

\begin{lstlisting}[language=C]
/*
 * TODO: by blocking fast commits here, we are increasing
 * fsync() latency slightly. Strictly speaking, we don't need
 * to block fast commits until the transaction enters T_FLUSH
 * state. So an optimization is possible where we block new fast
 * commits here and wait for existing ones to complete
 * just before we enter T_FLUSH. That way, the existing fast
 * commits and this full commit can proceed parallely.
 */
\end{lstlisting}

We moved the drain loop that waits for in-flight fast commits
from before \texttt{T\_LOCKED} to just before \texttt{T\_FLUSH},
along with the \texttt{journal->j\_fc\_off~=~0} reset that
depends on it. On our fio randwrite+fsync workload the patched
average commit time was 6661~$\mu$s versus 6590~$\mu$s on stock,
a 1\% delta in the wrong direction. The
\pth{/proc/fs/jbd2/loopX-8/info} counter showed the original
barrier almost never fired in this workload because
\texttt{T\_LOCKED} already takes well under a millisecond. An
earlier ``57\% improvement'' result turned out to be a setup
artifact: the stock benchmark had run for 5~MB while the patched
benchmark had run for 60 seconds, and the test filesystem was
not formatted with \texttt{-O fast\_commit} at all (so the drain
loop was never executed). After fixing both, no improvement
remained. Patch and postmortem are at
\pth{our_optimization/jbd2-fc-barrier-defer.patch} and
\pth{our_optimization/CANDIDATE1_postmortem.md}.

\subsection{C2: async bitmap prefetch completion}

The TODO sits at \texttt{fs/ext4/mballoc.c:2564}:

\begin{lstlisting}[language=C]
/*
 * TODO: We should actually kick off the buddy bitmap setup in a work
 * queue when the buffer I/O is completed, so that we don't block
 * waiting for the block allocation bitmap read to finish when
 * ext4_mb_prefetch_fini is called from ext4_mb_regular_allocator().
 */
\end{lstlisting}

We added a per-superblock workqueue and slab cache; in
\texttt{ext4\_mb\_prefetch\_fini}, if the bitmap is already
uptodate the original synchronous init is used, otherwise a work
item is queued and the caller returns. After two correctness
fixes caught by code review (a \texttt{noinline\_for\_stack}
attribute that we had silently stripped, and an OOM-path leak of
\texttt{s\_buddy\_cache}), the patch built cleanly. The problem
was again measurement: in our microbenchmark the cold-init path
was hit on $<$1\% of allocations (the bitmaps were already in
memory after the first few iterations), and our fio throughput
benchmark had a 70\% coefficient of variation across repeats on
the same kernel. We did not invest further. Patch and postmortem
are at \pth{our_optimization/mballoc-async-prefetch.patch} and
\pth{our_optimization/CANDIDATE2_postmortem.md}.

\paragraph{Lesson.} A patch is only as useful as the fraction of
operations that exercise the path it changes. After C1 and C2 we
adopted a hard rule: only ship a patch if a microbench can show
that the slow path it removes is hit on a substantial fraction of
the target operation.

\section{Candidate 3: Fast Commit Support for Inline xattrs}

\subsection{Intuition}

Every \texttt{setxattr} or \texttt{removexattr} call on a
fast-commit-enabled ext4 filesystem forces a full JBD2 commit.
The cost is paid on every SELinux relabel, every POSIX ACL
change, and every short \texttt{user.*} attribute write, which
makes the slow path effectively 100\% hit on xattr-touching
workloads. The ATC~2024 paper~\cite{atc24-fc} flags this gap but
provides no patch.

\subsection{Code analysis}

The fallback is fired by an unconditional call at the end of
\texttt{ext4\_xattr\_set\_handle} in \texttt{fs/ext4/xattr.c}
(line 2422):

\begin{lstlisting}[language=C]
ext4_fc_mark_ineligible(inode->i_sb, EXT4_FC_REASON_XATTR, handle);
\end{lstlisting}

A second \texttt{mark\_ineligible} call sits at line 2500 inside
the wrapper \texttt{ext4\_xattr\_set}, after
\texttt{ext4\_journal\_stop}, with \texttt{handle==NULL}. With
\texttt{handle} null this call marks the running JBD2 transaction
ineligible, which would suppress fast commit on any later
\texttt{fsync} sharing that transaction. Missing this second call
would have made the patch a no-op for end users, which is a key
finding from this candidate.

ext4 stores xattrs in three places: \emph{inline} (in the inode
body, between \texttt{EXT4\_GOOD\_OLD\_INODE\_SIZE
+ i\_extra\_isize} and the end of the inode), in a separate
\emph{block} (pointed to by \texttt{i\_file\_acl}), or in
dedicated \emph{EA inodes}. We targeted the inline case because
it covers nearly all real-world xattr operations: SELinux labels
and POSIX ACLs are short and fit inline.

\subsection{Design and implementation}

Two changes, both small.

\begin{enumerate}[leftmargin=*]
  \item \textbf{Expand the FC inode log.}
        \texttt{ext4\_fc\_write\_inode} in
        \texttt{fs/ext4/fast\_commit.c:863} currently logs the
        inode's bytes up to
        \texttt{EXT4\_GOOD\_OLD\_INODE\_SIZE
        + i\_extra\_isize}. We extend it to log up to
        \texttt{EXT4\_INODE\_SIZE} when the inode has inline
        xattrs (detected via the existing
        \texttt{EXT4\_STATE\_XATTR} flag). Replay needs no
        change: \texttt{ext4\_fc\_replay\_inode} already
        \texttt{memcpy}s \texttt{fc\_len} bytes, and the
        validator already accepts records up to
        \texttt{s\_inode\_size}.
  \item \textbf{Skip the ineligibility call when safe.} In
        \texttt{ext4\_xattr\_set\_handle}, add a local
        \texttt{bool touched\_block} that is set to true on every
        path that calls \texttt{ext4\_xattr\_block\_set} or sets
        \texttt{i.in\_inode = 1} (the EA-inode path). Replace the
        unconditional \texttt{mark\_ineligible} at line 2422 with:
\begin{lstlisting}[language=C]
if (error || touched_block || !ext4_has_feature_xattr(inode->i_sb))
    ext4_fc_mark_ineligible(inode->i_sb,
                            EXT4_FC_REASON_XATTR, handle);
\end{lstlisting}
        and delete the wrapper's redundant call at line 2500. The
        third condition handles the very first xattr on a fresh
        filesystem (the superblock xattr feature bit has just
        been set in memory but not yet persisted; if we
        fast-committed and the machine crashed, replay would not
        see the bit).
\end{enumerate}

The full patch is 108 lines (\pth{our_optimization/fc-inline-xattr.patch}):
9 added in \texttt{fast\_commit.c}, 7 added and 1 deleted in
\texttt{xattr.c}.

\subsection{Correctness}

\begin{itemize}[leftmargin=*]
  \item Three crash-recovery scripts
        (\texttt{c3\_crash\_test\_\{a,b,c\}.sh}). Test~A: 100
        inline xattrs, lazy unmount, remount, count
        (\textbf{PASS}, 100/100). Test~B: 100 set, 50 even
        removed, lazy unmount, remount, verify exactly the 50
        odd ones remain (\textbf{PASS}). Test~C: 500-byte xattr
        forces the block path, lazy unmount, remount, all 668
        bytes recovered (\textbf{PASS}, confirming the fallback
        branch still runs full commit when needed).
  \item xfstests subset
        \texttt{generic/\{062,118,300,337,454,455,473,482\}} +
        \texttt{ext4/032}. 5/5 runnable tests pass on both stock
        and patched; \texttt{generic/473} fails identically on
        both (a pre-existing issue) and three tests skip due to
        missing features. Logs at
        \pth{our_optimization/xfstests_c3/}.
\end{itemize}

\subsection{Results}

\paragraph{VM (VirtualBox, 11th Gen Intel Core i7-11800H, 4
cores, 5.7~GiB RAM; 256~MB loop fs with
\texttt{-O fast\_commit}; 5000
\texttt{fsetxattr}+\texttt{fsync} on a single inline-sized value).}

\begin{center}
\begin{tabular}{lrrr}
\toprule
Metric & Stock & Patched & Change \\
\midrule
Wall time         & 31{,}102~ms & 22{,}734~ms & $-26.9\%$ \\
JBD2 full commits & 5{,}000 & \textbf{78} & $\mathbf{-98.4\%}$ \\
\bottomrule
\end{tabular}
\end{center}

\paragraph{Bare-metal (11th Gen Intel Core i3-1115G4, 4
cores, 7.5~GiB RAM, real SATA SSD; 3 repeats).}

\begin{center}
\begin{tabular}{lrrr}
\toprule
Metric & Stock & Patched & Change \\
\midrule
Wall time (mean)  & 57{,}627~ms & 30{,}089~ms & $-47.8\%$ \\
JBD2 full commits & 5{,}000 & \textbf{78} & $\mathbf{-98.4\%}$ \\
\bottomrule
\end{tabular}
\end{center}

The transaction count reduction (5000 to 78, $\sim$64$\times$) is
identical between VM and bare-metal: this is the architectural
claim of the patch and it does not depend on the hardware. The
wall-time gain is larger on bare-metal (48\% versus 27\%) because
on a real SSD the fixed VFS and syscall overhead per
\texttt{fsync} is a smaller share of the total, leaving more room
for the journal-side savings to show through.

A non-xattr regression check (fio randwrite+fsync, 4 jobs, 30~s)
gave 528 IOPS patched versus 530 stock, well within noise.

\section{Candidate 4: Fast Commit Support for fallocate Range Operations}

\subsection{Intuition}

After C3 we asked which other \texttt{ext4\_fc\_mark\_ineligible}
fallback has the same shape: a 100\% hit rate on a well-defined
operation class, a small expected patch, and a microbench that
cleanly distinguishes patched from unpatched behavior.

\subsection{Characterization first}

For C4 we measured before writing kernel code. Three candidate
microbenchmarks were run on the C3 kernel: fallocate
COLLAPSE/INSERT range; large (4~KB) \texttt{setxattr} values that
spill to the xattr block (the case C3 does not cover); and
many-thread concurrent \texttt{fsync} (the CJFS-style compound
flush target). Three runs each, $N=1000$ ops (or 100 per thread):

\begin{center}
\begin{tabular}{lrrl}
\toprule
Workload & ms/op & tx/op & signal \\
\midrule
fallocate COLLAPSE\_RANGE & 7.03 & 1.001 & strong \\
fallocate INSERT\_RANGE   & 7.01 & 1.001 & strong \\
xattr block (4~KB value)  & 7.23 & 1.001 & strong but narrow \\
concurrent fsync $T=16$   & 1.21 & 0.003 & weak \\
\bottomrule
\end{tabular}
\end{center}

Fallocate COLLAPSE/INSERT showed the same per-op signal as C3's
xattr case. The xattr-block path showed equal per-op cost but
covers a minority of \texttt{setxattr} calls (typical SELinux/ACL
values are short and already handled by C3). Concurrent fsync was
already well-coalesced by JBD2's group commit
(tx/op falls from 0.020 at $T=1$ to 0.003 at $T=16$), so a
CJFS-style port would not help much on 6.1.4. Fallocate became
C4.

\subsection{Code analysis}

Two call sites in \texttt{fs/ext4/extents.c} mark fallocate
range operations as ineligible, both right after
\texttt{ext4\_journal\_start}:

\begin{lstlisting}[language=C]
/* line 5363, inside ext4_collapse_range */
ext4_fc_mark_ineligible(sb, EXT4_FC_REASON_FALLOC_RANGE, handle);

/* line 5508, inside ext4_insert_range */
ext4_fc_mark_ineligible(sb, EXT4_FC_REASON_FALLOC_RANGE, handle);
\end{lstlisting}

Crucially, ext4 already has fast-commit tags for the underlying
extent-level changes
(\texttt{FC\_TAG\_ADD\_RANGE}, \texttt{FC\_TAG\_DEL\_RANGE},
emitted by \texttt{ext4\_fc\_write\_inode\_data}), and the
per-inode \texttt{FC\_TAG\_INODE} tag carries the new
\texttt{i\_size}. The fallback exists only because the
collapse/insert paths never call
\texttt{ext4\_fc\_track\_range} on the logical blocks they
modify, so the FC commit walker does not know to include the
shifted region.

\subsection{Design and implementation}

Replace each of the two \texttt{mark\_ineligible} calls with a
call to \texttt{ext4\_fc\_track\_range} over the affected
logical-block range, computed at the point where the journal
handle is open:

\begin{lstlisting}[language=C]
/* collapse path: track punch_start to current EOF */
ext4_fc_track_range(handle, inode, punch_start,
    ((inode->i_size + inode->i_sb->s_blocksize - 1) >>
     EXT4_BLOCK_SIZE_BITS(inode->i_sb)) - 1);

/* insert path: track offset_lblk to post-grow EOF */
ext4_fc_track_range(handle, inode, offset_lblk,
    ((inode->i_size + len + inode->i_sb->s_blocksize - 1) >>
     EXT4_BLOCK_SIZE_BITS(inode->i_sb)) - 1);
\end{lstlisting}

At fsync time, \texttt{ext4\_fc\_write\_inode\_data} walks the
post-shift extent tree over the tracked range and emits
\texttt{FC\_TAG\_ADD\_RANGE} for each shifted extent at its new
position plus \texttt{FC\_TAG\_DEL\_RANGE} for any resulting
hole. \texttt{FC\_TAG\_INODE} is emitted by the
\texttt{ext4\_mark\_inode\_dirty} call later in each function
and carries the new \texttt{i\_size}.

\paragraph{Replay correctness.} The non-obvious part is that
\texttt{ext4\_fc\_replay\_del\_range} frees the physical blocks
of the range it removes via \texttt{ext4\_mb\_mark\_bb(\ldots,
0)}. In an insert (where the physical data does not move, only
the logical position shifts) \texttt{DEL\_RANGE} replay
transiently marks blocks free and \texttt{ADD\_RANGE} replay
re-binds them. The intermediate bitmap state is wrong; the final
state is reconciled by the pre-existing
\texttt{ext4\_fc\_set\_bitmaps\_and\_counters} pass, which runs
at the end of replay and re-marks all currently-reachable blocks
as used.

The full patch is 46 lines (\pth{our_optimization/fc-fallocate-range.patch}),
two hunks in \texttt{fs/ext4/extents.c}. No change to
\texttt{fast\_commit.c}, no new tag types, no on-disk format
change.

\subsection{Correctness}

\begin{itemize}[leftmargin=*]
  \item Three crash-recovery scripts
        (\texttt{c4\_crash\_test\_\{a,b,c\}.sh}). Each writes a
        distinct-byte-pattern file, performs the operation,
        fsyncs, lazy-umounts to simulate a crash, remounts, and
        compares the resulting file's md5 and size to the
        expected post-operation state computed in pure Python.
        Test~A: collapse 32 blocks (\textbf{PASS}). Test~B:
        insert 16 blocks at offset~16 (\textbf{PASS}). Test~C:
        three interleaved collapse/insert ops (\textbf{PASS}).
        Test~C is the strongest of the three because it stresses
        replay with multiple ADD\_RANGE and DEL\_RANGE tags
        accumulated for the same inode in a single FC log.
  \item Same xfstests subset as C3. 5/5 runnable tests pass;
        \texttt{generic/473} fails identically to stock and to
        the C3 kernel. \texttt{generic/062, 337, 454} (the
        xattr-heavy tests) all pass, confirming C4 does not
        regress C3's work. Logs at
        \pth{our_optimization/xfstests_c4/patched_6.1.4-cs614-c4-patched.log}.
\end{itemize}

\subsection{Results}

\paragraph{VM (same setup as C3: i7-11800H, 4 cores, 5.7~GiB RAM,
256~MB loop fs with \texttt{-O fast\_commit}; 1000 ops + per-op
fsync, three runs per mode).} The C3 kernel is the baseline
because C3 does not change any fallocate code; on this path the
C3 kernel and stock 6.1.4 are architecturally identical, and our
characterization runs confirmed tx/op~=~1.001 on both.

\begin{center}
\begin{tabular}{lrrr}
\toprule
Metric & Baseline (C3) & Patched (C4) & Change \\
\midrule
\multicolumn{4}{l}{\textit{COLLAPSE\_RANGE}} \\
Wall time (mean)  & 7{,}029~ms & 5{,}423~ms & $-22.8\%$ \\
JBD2 full commits & 1{,}001 & \textbf{16} & $\mathbf{-98.4\%}$ \\
\midrule
\multicolumn{4}{l}{\textit{INSERT\_RANGE}} \\
Wall time (mean)  & 7{,}006~ms & 5{,}189~ms & $-25.9\%$ \\
JBD2 full commits & 1{,}001 & \textbf{16} & $\mathbf{-98.4\%}$ \\
\bottomrule
\end{tabular}
\end{center}

\paragraph{Bare-metal (11th Gen Intel Core i3-1115G4, 4
cores, 7.5~GiB RAM, real SATA SSD; 3 repeats).}

\begin{center}
\begin{tabular}{lrrr}
\toprule
Metric & Baseline (C3) & Patched (C3+C4) & Change \\
\midrule
\multicolumn{4}{l}{\textit{COLLAPSE\_RANGE}} \\
Wall time (mean)  & 12{,}445~ms & 7{,}032~ms & $-43.5\%$ \\
JBD2 full commits & 1{,}001 & \textbf{16} & $\mathbf{-98.4\%}$ \\
\midrule
\multicolumn{4}{l}{\textit{INSERT\_RANGE}} \\
Wall time (mean)  & 11{,}539~ms & 6{,}680~ms & $-42.1\%$ \\
JBD2 full commits & 1{,}001 & \textbf{16} & $\mathbf{-98.4\%}$ \\
\bottomrule
\end{tabular}
\end{center}

The transaction-count reduction is byte-identical to the VM
(1001~$\to$~16 in both modes). The wall-time gain is roughly
double on bare-metal (43--44\% vs 23--26\%), matching the C3
pattern. Patched-run variance is small: 0.5\% CV for COLLAPSE,
2\% for INSERT.

\section{Discussion}

\subsection{Why C3 and C4 are measurable and C1, C2 are not}

\begin{center}
\begin{tabular}{lcccc}
\toprule
 & C1 & C2 & \textbf{C3} & \textbf{C4} \\
\midrule
Slow-path hit rate & $\sim$0.1\% & $<$1\% & \textbf{100\%} & \textbf{100\%} \\
Effect per hit     & $<$1~ms     & sub-ms & 5--11~ms        & 5--7~ms \\
Measurable on VM?  & No          & No     & Yes (64$\times$) & Yes (62.5$\times$) \\
\bottomrule
\end{tabular}
\end{center}

C3 and C4 together remove two of the fast-commit fallback reasons
documented in \texttt{fs/ext4/fast\_commit.h}: \texttt{XATTR}
(inline case) and \texttt{FALLOC\_RANGE} (collapse and insert).
The two patches touch disjoint files (\texttt{xattr.c} and
\texttt{fast\_commit.c} for C3; \texttt{extents.c} for C4), so
they compose without conflict.

\subsection{Why wall-time gain is smaller than commit-count gain}

For both patches, commit-count drops by $\sim$60$\times$ but
wall-time drops by only $\sim$1.4$\times$ on the VM and
$\sim$2$\times$ on bare-metal. The gap is the fixed per-fsync
overhead our patches do not touch: VFS and syscall entry/exit,
device write-behind, and (for C4) extent-tree manipulation in
\texttt{ext4\_ext\_remove\_space} and
\texttt{ext4\_ext\_shift\_extents}. Bare-metal closes the gap
because real SSD has a smaller share of non-journal cost per
\texttt{fsync}.

\subsection{Scope limits}

We did not extend C3 to xattrs that overflow into a separate block
or into EA inodes. Doing so would require new fast-commit tag
types (\texttt{FC\_TAG\_XATTR\_SET}, \texttt{FC\_TAG\_XATTR\_DEL})
and replay support, which in turn would need a matching
e2fsprogs patch and an on-disk compat feature bit. Our
characterization confirms the per-op signal is as strong as the
inline case, but the real-world hit rate is much smaller because
typical SELinux/ACL values are short. We did not address
\texttt{CROSS\_RENAME} or \texttt{RENAME\_DIR} either, since
their fallback exists for a deeper reason: replay cannot yet
update the \texttt{..} dirent safely without a non-trivial
rewrite.

\section{Conclusion}

We submit two small, independent patches against ext4 on
Linux~6.1.4 that close two fast-commit fallback reasons
(\texttt{XATTR} inline-case and \texttt{FALLOC\_RANGE}). Both
produce a $\sim$60$\times$ reduction in JBD2 full commits on
their respective microbench, both reproduce that
architectural claim byte-identically on bare-metal, and both
preserve crash-recovery semantics under custom and xfstests-based
testing. Two earlier candidates were correctly implemented but
optimized paths too rarely hit to measure; their postmortems are
included in the artifact for completeness.

\begin{thebibliography}{9}
\bibitem{ext4-design}
A. Mathur, M. Cao, S. Bhattacharya, A. Dilger, A. Tomas, L. Vivier.
\textit{The new ext4 filesystem: current status and future plans}.
Proceedings of the Linux Symposium, Ottawa, 2007.
\url{https://www.kernel.org/doc/ols/2007/ols2007v2-pages-21-34.pdf}

\bibitem{ijournaling}
D. Park, D. Shin.
\textit{iJournaling: Fine-Grained Journaling for Improving the
Latency of Fsync System Call}.
USENIX ATC 2017.
\url{https://www.usenix.org/conference/atc17/technical-sessions/presentation/park}

\bibitem{optfs}
V. Chidambaram, T. S. Pillai, A. C. Arpaci-Dusseau, R. H.
Arpaci-Dusseau.
\textit{Optimistic Crash Consistency}.
ACM SOSP 2013.
\url{https://research.cs.wisc.edu/adsl/Publications/optfs-sosp13.pdf}

\bibitem{barrierfs}
Y. Won, J. Jung, G. Choi, J. Oh, S. Son, J. Hwang, S. Cho.
\textit{Barrier-Enabled IO Stack for Flash Storage}.
USENIX FAST 2018.
\url{https://www.usenix.org/conference/fast18/presentation/won}

\bibitem{atc24-fc}
H. Shirwadkar, S. Kadekodi, T. Ts'o.
\textit{FastCommit: Resource-Efficient, Performant, and
Cost-Effective File System Journaling}.
USENIX ATC 2024.

\bibitem{cjfs}
Y. Oh, J. Yoo, D. Nam, C. Min, Y. Won.
\textit{CJFS: Concurrent Journaling for Better Scalability}.
USENIX FAST 2023.

\bibitem{djfs}
J. Yoo, Y. Oh, et al.
\textit{DJFS: Directory-Granularity Filesystem Journaling for
CMM-H SSDs}.
USENIX FAST 2025.

\bibitem{syncsync}
Q. Jiang, et al.
\textit{Sync+Sync: A Covert Channel Built on fsync with Storage}.
USENIX Security 2024.
\end{thebibliography}

\end{document}
